% Figure — Oscillateur Colpitts à source commune (a) et schéma équivalent simplifié, sortie en Thévenin (b) (EANA3, chapitre 4).
% D'après les notes externes [notes-2023-c03-05, p. 3-4], redessinée (polarisation non montrée, comme dans les notes).
% Rendu SVG :
%   latex -output-directory=<dossier temporaire> 04-colpitts.tex
%   dvisvgm --no-fonts --exact-bbox -o 04-colpitts.svg <dossier temporaire>/04-colpitts.dvi
\documentclass[border=4pt]{standalone}
\usepackage[T1]{fontenc}
\usepackage[utf8]{inputenc}
\usepackage{amsmath}
\usepackage{sansmath}
\renewcommand{\familydefault}{\sfdefault}
\usepackage[european, straightvoltages, americancurrents]{circuitikz}
\usetikzlibrary{shapes.geometric,calc,positioning}
\begin{document}
\sansmath
\begin{circuitikz}[font=\small, >=latex, fleche/.style={->, thick}]
  % (a)
  \node at (-0.4,4.4) {(a)};
  \draw (2.6,1.6) node[nmos] (M) {};  \node[right] at (M.text) {$M$};
  \draw (M.S) node[ground]{};
  \draw (M.D) -- ++(0,0.4) coordinate (d) to[isource, invert, l=$I_{bias}$] ++(0,1.8) node[vcc]{$V_{DD}$};
  \draw (d) -- (4.6,0 |- d) coordinate (dr) to[C, l=$C_1$] (4.6,0) node[ground]{};
  \draw (M.G) -- (0.6,0 |- M.G) coordinate (g) to[C, l_=$C_2$] (0.6,0) node[ground]{};
  \draw (g) -- (0,0 |- g) -- (0,-0.9) to[cute inductor, l_=$L$] (5.3,-0.9) -- (5.3,0 |- dr) -- (dr);  \fill (g) circle (1.5pt);  \fill (dr) circle (1.5pt);
  % (b)
  \begin{scope}[xshift=8.2cm]
    \node at (-0.4,4.4) {(b)};
    \draw (0,2.6) to[american controlled voltage source, l_=$-g_m r_{ds}\, v_{gs}$] (0,0);  \draw (0,2.6) to[R, l=$r_{ds}$, i>^=$i_1$] (2.4,2.6) to[cute inductor, l=$L$, i>^=$i_2$] (5.0,2.6);
    \draw (2.4,2.6) to[C, l=$C_1'$] (2.4,0);  \draw (5.0,2.6) to[C, l=$C_2'$] (5.0,0);
    \draw (0,0) -- (5.0,0);  \draw (2.4,0) node[ground]{};
    \draw[fleche] (6.3,0.2) -- node[right]{$v_{gs}$} (6.3,2.4);  \draw (5.0,2.6) -- (6.3,2.6);  \draw (5.0,0) -- (6.3,0);
  \end{scope}
\end{circuitikz}
\end{document}
